%% Chapter 3 - Research Methodology
%% Planned sections (master structure): 3.1 Research design; 3.2 Data participants
%% or study materials; 3.3 Experimental or study procedure; 3.4 Measures and
%% evaluation metrics; 3.5 Data analysis; 3.6 Ethics privacy security and
%% responsible AI; 3.7 Reproducibility and research artefacts.
%% All section labels are now defined: sec:research-design (3.1),
%% sec:data-materials (3.2), sec:procedure (3.3), sec:measures (3.4),
%% sec:data-analysis (3.5), sec:ethics (3.6), sec:reproducibility (3.7).
%% CANONICAL RESULTS SOURCE (author-confirmed, 2026-10-06): the 1500-record
%% evaluation in "Thesis/Reserch Paper/NeuroShell_IRE - Reserch paper.tex"
%% (N=1500, 600 well-formed + 900 flawed; Exp1--Exp4 tables; per-backend
%% ablation; single-run point estimates, no confidence intervals). The 200-record
%% run at git HEAD was a SMOKE TEST only. Latency figures are NOT in the 1500
%% draft; the final paper ("IRE-Doc/Research Paper/..._Research_Paper.tex")
%% reports p50 16.5 s, p95 70.9 s, cache-hit 44 ms on the same hardware --
%% reconcile with the author before quoting.
\chapter{Research Methodology}\label{ch:methodology}

\section{Research Design}\label{sec:research-design}

This chapter defines how the research questions posed in
Section~\ref{sec:research-questions} are answered through measurement. The
dissertation is a design-and-evaluation study: it designs and implements an
engineered artefact---the guarded intent-recognition engine introduced in
Section~\ref{sec:overall-project-context}---and evaluates that artefact
empirically under controlled, reproducible conditions. The artefact is the
object under study, and the controlled experiment over an annotated data set is
the measurement instrument. The design is committed in advance: a fixed protocol
applied to a fixed data set with a versioned evaluation harness
(Section~\ref{sec:reproducibility}). The study is deliberately not a
hypothesis-testing one; its outcomes are quantitative point estimates from a
committed experiment record, and no statistical significance claim is attached
to them (Sections~\ref{sec:research-questions} and
\ref{sec:scope-limitations-assumptions}).

\subsection*{System under study}

The object evaluated is the guarded pipeline. At the level required to fix the
measurement design, the pipeline is a seven-stage core that converts a free-form
utterance into a validated intent---normalisation, alias resolution, model
inference, parsing, schema validation, regex and entity validation, and scope
enforcement---wrapped in five middleware components (the adversarial detector,
pre-inference role check, session context, the semantic cache, and the
post-validation role check) under a two-sided role-based access-control (RBAC)
policy. The measurement design needs only the guard predicates: their presence
or absence defines the treatment conditions. The complete stage-by-stage
definition and interface contract are developed in
Chapter~\ref{ch:product-design}.

Following the framework of Section~\ref{sec:gap-synthesis}, let $\mathcal{X}$
denote the utterance space, $\mathcal{Y}$ the closed intent taxonomy,
$\mathcal{C}$ the set of classified flaw categories and $\rho$ the caller role.
The decoder is a stochastic map $f:\mathcal{X}\to\mathcal{Z}$; the deterministic
guards are the predicates $g_5,g_6,g_{6B},g_7$ over the schema, regex/entity,
network and scope stages, the RBAC authorisation $b:\mathcal{R}\times\mathcal{Y}
\to\{0,1\}$, and $m_1$ the adversarial-input detector. The guarded decision
function is
\begin{equation}
\Phi(x)=\begin{cases}
y\in\mathcal{Y}\ \text{accepted}, & \text{if } G(z,x)\wedge b(\rho,y)\wedge\neg\, m_1(x),\\
e\in\mathcal{C}\ \text{rejected}, & \text{otherwise},
\end{cases}
\label{eq:guarded-decision}
\end{equation}
where $G(z,x)=g_5(z)\wedge g_6(z)\wedge g_{6B}(z,x)\wedge g_7(x)$. The safety
invariant that the design enforces is that every accepted intent satisfies every
security predicate:
\begin{equation}
\mathrm{Accepted}\!\big(\Phi(x)\big)\;\Rightarrow\;\bigwedge\nolimits_{k\in K}
P_k(x),\qquad K=\{\mathrm{S5},\mathrm{S6},\mathrm{S6B},\mathrm{S7},
\mathrm{RBAC},\mathrm{M1}\}.
\label{eq:safety-invariant}
\end{equation}
Experiment~1 operationalises RQ1 as the gap in
$\Pr[\bigwedge_{k\in K}P_k]$ between the raw decoder $f$ and the guarded
decision $\Phi$; Experiment~3 operationalises RQ2 as the change in the catch
rate attributable to each predicate when that predicate is removed from the
conjunction (Section~\ref{sec:procedure}).

\subsection*{Treatment structure}

The measurement design manipulates exactly three factors, and nothing else is
allowed to vary between conditions:
\begin{itemize}
\item \textbf{Pipeline condition} (RQ1; Experiment~1). Raw single-shot model
output is compared with the fully guarded pipeline, holding the backend, data set
and thresholds fixed.
\item \textbf{Backend} (RQ3; Experiments~2 and~3). Two offline,
instruction-tuned backends---\texttt{qwen2.5-coder:7b} and
\texttt{gemma4:e4b}---are compared under the identical guardrail stack and
identical prompts.
\item \textbf{Configuration} (RQ2; Experiment~3). The fully guarded baseline is
compared with seven ablation conditions, each removing exactly one guard
component: the scope guard, RBAC, the network validator, the hallucination
validator, the deterministic validation stack as a combined block, the
self-consistency gate, and the semantic cache. The set of ablation conditions is
fixed in advance so that each layer's marginal contribution is measured against a
constant baseline.
\end{itemize}

The dependent variables are the three rates that measure whether the system
``does the right thing and does not do the wrong thing'' under a guard: intent
accuracy over the well-formed set, the flaw-catch rate over the flawed set, and
the guard false-positive rate over the well-formed set. Their formal definitions
and the strictness of the flaw-catch criterion are given in
Section~\ref{sec:measures}. A fourth experiment (Experiment~4) measures catch
separately for each flaw category, so that the aggregate ablation result is
interpreted per class rather than through a single number.

\subsection*{Control variables}

To keep every comparison interpretable, the following are held constant across
all conditions: the evaluation data set (Section~\ref{sec:data-materials}); the
hardware platform and the CPU-only Ollama runtime; the sampling temperatures
(0.4 for the primary inference and 0.6 for the $k=3$ consistency samples); the
semantic agreement threshold; and the operator-configured engagement scope and
RBAC policy under evaluation. Only the manipulated factor differs between
conditions.

\subsection*{Measurement epistemology}

Deterministic guard layers produce an exact outcome for a fixed input; the
stochastic model component is executed once per input and condition. Every
figure reported in this dissertation is therefore a single-run point estimate
without seed sweeps or confidence intervals. The design compensates by
committing the complete protocol, data set and harness to a versioned evaluation
manifest, so that the record can be reproduced, audited and extended
(Section~\ref{sec:reproducibility}). This choice is consistent with the
evidential boundary declared in Section~\ref{sec:scope-limitations-assumptions},
and its consequences for how the results may be read are discussed in
Chapter~\ref{ch:discussion}.

\subsection*{Chapter organisation}

The remainder of this chapter defines the study materials and the constitution
of the data set (Section~\ref{sec:data-materials}), the experimental procedure
and ablation protocol (Section~\ref{sec:procedure}), the metrics
(Section~\ref{sec:measures}), the method of analysis applied to the recorded
measurements (Section~\ref{sec:data-analysis}), the ethics, privacy, security
and responsible-AI considerations (Section~\ref{sec:ethics}), and the
reproducibility and research-artefact commitments (Section~\ref{sec:reproducibility}).
The implementation of the system under test is described in
Chapter~\ref{ch:implementation-testing}, the recorded outcomes in
Chapter~\ref{ch:results}, and their interpretation against RQ1--RQ3 in
Chapter~\ref{ch:discussion}.

\TODO{At the final AUDIT step, confirm the holdout definition used for
Experiments~1 and~2 against the committed experiment record: the canonical
1500-record draft states that the stratified full 1500-prompt set was used,
whereas the older 200-record paper draft reports a 200-prompt holdout. The
chapters must quote whichever definition the versioned evaluation manifest
records.}

\section{Data and Study Materials}\label{sec:data-materials}

This study involves no human participants. The study materials are the annotated
evaluation data set that drives the recorded experiments, the deployment
platform on which they run, and the two language-model backends under test. The
data set is author-authored rather than sampled from operational traffic, an
assumption declared in Section~\ref{sec:scope-limitations-assumptions}; it is
versioned with the project source and regenerable from the synthetic-generation
script used to author it (Section~\ref{sec:reproducibility}).

\subsection*{Constitution of the data set}

The evaluation data set $D$ has two disjoint subsets. The well-formed subset
$D^+$ contains 600 free-form security requests, distributed across the nine
top-level intents of the closed taxonomy used by the system under test
(scanning and discovery, exploitation, password attack, directory bruteforce,
CVE lookup, reconnaissance, phishing, network defence, and miscellaneous).
Each record in $D^+$ is annotated with a gold intent label $y^\ast(x)$.

The flawed subset $D^-$ contains 900 deliberately malfunctioning requests,
distributed across the eight flaw categories in Table~\ref{tab:dataset-composition}.
Each record in $D^-$ carries two gold labels: an intent label $y^\ast(x)$ and a
flaw label $c^\ast(x)$ drawn from the eight-category flaw taxonomy $\mathcal{C}$,
which is the rejection category a correct pipeline is expected to assign
(Section~\ref{sec:measures}). The eight categories span the failure modes that
matter in an automated security toolchain: adversarial phraseology intended to
divert the model; requests whose target or technique falls outside the
operator-declared engagement scope; requests too ambiguous or under-specified to
resolve; internally contradictory requests; fabricated parameter values;
fabricated CVE identifiers; requests whose phrasing mismatches the structure of
the stated intent; and requests an unauthorised role attempts to issue. The
mapping between each category and the guard layer or layers responsible for
catching it is a hypothesis tested by the per-class experiment
(Section~\ref{sec:procedure}) and reported in Chapter~\ref{ch:results}.

\begin{table}[h]
\centering
\footnotesize
\caption{Experimental data set composition ($N=1500$). $D^+$ holds 600
well-formed requests across the nine-intent taxonomy; $D^-$ holds 900
deliberately flawed requests across eight flaw categories.}
\label{tab:dataset-composition}
\begin{tabular}{@{}lc@{}}
\hline
\textbf{Category} & \textbf{Records} \\
\hline
$D^+$ well-formed requests & 600 \\
Adversarial injection & 113 \\
Out-of-scope target/technique & 113 \\
Ambiguous / under-specified & 113 \\
Contradictory request & 113 \\
Fabricated parameter & 112 \\
Fabricated CVE & 112 \\
Type mismatch (string vs. intent) & 112 \\
RBAC violation & 112 \\
\hline
Total & 1500 \\
\hline
\end{tabular}
\end{table}

The flawed categories in the 1500-record set are near-evenly sized (113 or 112
records each, the $D^-$ classes summing to 900). Per-class catch rates are
therefore reported against each category's own denominator rather than as a
single aggregate, and cross-category comparisons are only meaningful against
those denominators (Section~\ref{sec:measures}). This is also why the per-class
experiment supplements the aggregate ablation: no flawed category is large
enough on its own to dominate the aggregate figure.

\subsection*{Annotation practice}

Gold labels are assigned at authoring time under a fixed category definition, so
that every record in a category is flawed for the stated reason and no record is
flawed for a reason that belongs to another category. Because the flaw-catch
metric requires the assigned rejection category to match the gold label exactly
(Section~\ref{sec:measures}), borderline phrasings that sit between two
categories contribute label noise, most plausibly at the boundary between the
adjacent semantic classes (ambiguous and contradictory); the residual weakness
of those classes is reported in Chapter~\ref{ch:results} and their
interpretation is qualified accordingly.

\subsection*{Platform and backends}

All recorded measurements were obtained on a single consumer laptop---an ASUS
VivoBook X513EA (Intel Core i5-1135G7, four cores and eight threads, 24~GB
DDR4, Windows 11)---running Ollama in CPU-only mode; no GPU was used. The two
backends under test are \texttt{qwen2.5-coder:7b} (4.7~GB on disk) and
\texttt{gemma4:e4b} (9.6~GB on disk), the build to which the
\texttt{gemma4:latest} alias resolved at evaluation time. The Ollama version and
the per-model image digests as fetched are pinned in the evaluation manifest,
which is the authoritative anchor for reproducing the study
(Section~\ref{sec:reproducibility}). The engagement scope, the RBAC policy and
the semantic threshold used by the self-consistency gate are configuration
materials of the same manifest; they are held constant across all conditions
(Section~\ref{sec:research-design}).

\TODO{At the final AUDIT step, reconcile the composition reported here
($N=1500$; $D^+{=}600$; $D^-{=}900$ with the eight flawed categories) against the
committed synthetic dataset generation script and its generated artefacts, and
confirm the annotation conventions against the golden dataset stored with the
project source.}

\section{Experimental Procedure}\label{sec:procedure}

This section defines how the materials of Section~\ref{sec:data-materials} are
processed into the recorded measurements reported in
Chapter~\ref{ch:results}. The procedure is fixed in advance and versioned: every
pass is executed by the evaluation harness against the same data set in the same
record order, and the decision recorded for every record is written to the
committed experiment record (Section~\ref{sec:reproducibility}).

\subsection*{Execution model}

For each record, the harness issues the request to the service once per
condition, and buffers the single recorded decision. In the \emph{guarded}
condition the pipeline runs end-to-end: the pre-inference role check, the
adversarial detector, the seven core stages including the self-consistency gate,
and the post-validation role check. The outcome is either an accepted intent
$y\in\mathcal{Y}$ (only when the conjunction of guards holds,
Equation~\ref{eq:guarded-decision}) or a rejection $e\in\mathcal{C}$ carrying
the guard-assigned flaw category. In the \emph{raw} condition, used by
Experiment~1 only, the harness records the model's single-shot structured output
and applies no guard predicates: no rejection category exists in raw output, so
no raw output can be classified as flawed.

\subsection*{Experiment set}

\begin{itemize}
\item \textbf{Experiment 1 (pipeline value; RQ1).} For each backend, the raw
condition is compared with the guarded condition over the full stratified
1500-prompt set of Section~\ref{sec:data-materials} (600 well-formed and 900
flawed records). This quantifies the contribution of the guardrail stack as the
gap in the probability that output satisfies the full conjunction of security
predicates (Section~\ref{sec:research-design}).
\item \textbf{Experiment 2 (backend head-to-head; RQ3).} The guarded stack is
held constant while the two backends are compared on the identical full
holdout with identical prompts. This isolates the effect of the underlying model
from every other pipeline component.
\item \textbf{Experiment 3 (layer ablation; RQ2 and RQ3).} Both backends are
evaluated on the full 1500-record set under eight configurations: the fully
guarded baseline and seven ablation conditions, each removing exactly one guard
component---the scope guard, RBAC, the network validator, the hallucination
validator, the combined deterministic validation stack (schema, regex/entity,
network and scope validation disabled together as one block), the
self-consistency gate, and the semantic cache. The combined validation-stack
condition is a \emph{joint} ablation that bounds the stack's aggregate
contribution; it is not a single layer. The marginal contribution of a layer is
\begin{equation}
\Delta\mathrm{Catch}_{\ell}=\mathrm{Catch}_{\Pi\setminus\ell}
-\mathrm{Catch}_{\Pi},
\label{eq:delta-catch}
\end{equation}
so that a negative value indicates the layer contributes to the catch rate.
\item \textbf{Experiment 4 (per-class catch; RQ2 evidence).} Catch is measured
separately for each of the eight flaw categories on a dedicated deep-dive subset
of flawed prompts, so that the aggregate ablation result is interpreted against
each category's own denominator rather than a single figure.
\end{itemize}

\subsection*{Self-consistency gate}

The self-consistency gate adapts the sampling-based consistency strategy of
Wang \emph{et al.} from generative reasoning to structured classification: for
each primary inference (temperature $0.4$), $k=3$ additional responses are
sampled at temperature $0.6$, embedded, and compared for agreement. The recorded
agreement signal is
\begin{equation}
\alpha(x)=\binom{k}{2}^{-1}\sum_{1\le i<j\le k}
\mathbf{1}\big[d_{ij}\le\theta_s\big],
\label{eq:agreement-signal}
\end{equation}
where $d_{ij}=1-\cos\big(\mathbf{e}(z_i),\mathbf{e}(z_j)\big)$, $\mathbf{e}(\cdot)$
is the sentence embedding, and $\theta_s$ is the semantic threshold pinned in the
evaluation manifest; a raw single sample has $\alpha=1$ by construction.
Agreement below the threshold demotes confidence and routes the request to the
hallucination validator, which assigns the output to one of six error classes
(impossible CVE, impossible parameter range, unverifiable value, contradictory
fields, out-of-scope reference, unresolved type) and selects the corresponding
response---\textsc{block}, \textsc{warn}, or \textsc{escalate} (the policy is a
configuration of the manifest and is exercised by the automated test suite).

\subsection*{Adversarial detection}

The adversarial detector precedes the model where possible. Its lexical layer
assigns suspicion weights of $0.8$--$1.0$ to prompt tokens, its structural layer
scores sentence topology and imperative-but-unclassifiable phrasing at
$0.3$--$0.9$, and the scoring layer fuses the two signals cumulatively. A
cumulative score above $0.7$ flags the request for heightened scrutiny, and a
score of at least $0.95$ auto-blocks it without paying for model inference.

\subsection*{Fixed operational parameters}

The following are held constant across every condition (Section~\ref{sec:research-design})
and pinned in the evaluation manifest: the primary sampling temperature ($0.4$);
the consistency sampling count ($k=3$) and temperature ($0.6$); the semantic
agreement threshold $\theta_s$; the adversarial-detector flag and auto-block
thresholds ($0.7$ and $0.95$); the operator-configured engagement scope; the RBAC
policy; and a per-stage timeout so that an unresponsive model cannot stall the
pipeline. Design decisions that would add a further model call were rejected on
latency grounds, which bounds the sampling budget at the fixed $k$ above.

\subsection*{Recording and reproduction}

Each configuration is executed once per backend, with the stochastic model
component sampled once per record and condition (Section~\ref{sec:research-design}).
Per-record decisions are appended to the experiment record in manifest-defined
order, and every table reported in Chapter~\ref{ch:results} is regenerated from
that record by a single command (Section~\ref{sec:reproducibility}), so the
measured statements of this dissertation are, in principle, reproducible without
re-running the model.

\TODO{At the final AUDIT step, (i) confirm the composition of the
Experiment~4 deep-dive flawed subset and its per-class denominators against the
committed manifest, since the canonical 1500-record draft does not enumerate
them, and (ii) re-verify how the raw-condition false-positive figures were
computed by the evaluation script given that the metrics definition reserves
false positives for guard flags (Section~\ref{sec:measures}).}

\section{Measures and Evaluation Metrics}\label{sec:measures}

Three rates constitute the dependent variables of the study
(Section~\ref{sec:research-design}). Let $y^\ast(x)$ and $c^\ast(x)$ denote the
gold intent and flaw labels of record $x$, and $\hat{y}(x)$ and $\hat{c}(x)$ the
intent and flaw category assigned by the system. In the guarded condition an
accepted record carries a recognised intent $\hat{y}(x)$ and no flaw category;
a rejected record carries a classified flaw category $\hat{c}(x)$ and no intent.
Two conventions follow: a well-formed record that a guard flags is counted as
not recognised (so it never inflates accuracy), and a flawed record that is
rejected under \emph{any} other category is counted as not caught. All rates are
reported as percentages.

\subsection*{Intent accuracy}

The intent accuracy $D_{\mathrm{Acc}}$ is computed over the legitimate-input set
$D^+$:
\begin{equation}
D_{\mathrm{Acc}}=\frac{100}{|D^+|}\sum_{x\in D^+}
\mathbf{1}\!\left[\hat{y}(x)=y^\ast(x)\right],
\label{eq:metric-dacc}
\end{equation}
where $\mathbf{1}[\cdot]$ is the indicator function. It measures how often the
pipeline both accepts a well-formed request \emph{and} recovers its intended
class---the system ``doing the right thing.''

\subsection*{Flaw-catch rate}

The flaw-catch rate is computed over the flawed-input set $D^-$:
\begin{equation}
\mathrm{Catch}=\frac{100}{|D^-|}\sum_{x\in D^-}
\mathbf{1}\!\left[\hat{c}(x)=c^\ast(x)\right].
\label{eq:metric-catch}
\end{equation}
The strictness of this criterion matters for interpretation: a request that is
rejected, but under the wrong category, is a miss. The metric therefore rewards
not merely refusal but \emph{correctly diagnosed} refusal, which is the property
an automated toolchain needs---a downstream consumer must know \emph{why} a
request was refused, not only that it was. Aggregate catch is complemented by
per-class reporting (Experiment~4) because the flaw categories of the data set
carry their own denominators and no single figure expresses the coverage of
every category (Section~\ref{sec:data-materials}).

\subsection*{False-positive rate}

The false-positive rate is computed over the legitimate-input set $D^+$:
\begin{equation}
\mathrm{FP}=\frac{100}{|D^+|}\sum_{x\in D^+}
\mathbf{1}\!\left[\text{a guard flags }x\right],
\label{eq:metric-fp}
\end{equation}
measuring how often a guard erroneously flags a well-formed request (a
``shadow-ban''). A flagged well-formed request is a conservative over-block, not
a misclassification toward a privileged intent: the guard either marks the
request for heightened scrutiny or rejects it outright; it never re-approves it
as a different intent. FP is therefore the usability counterpart of the two
safety rates, and the design goal is simultaneously high $D_{\mathrm{Acc}}$ and
high Catch with low FP.

\subsection*{Ablation measure}

The per-layer contribution to catch is the delta defined in
Section~\ref{sec:procedure},
$\Delta\mathrm{Catch}_{\ell}=\mathrm{Catch}_{\Pi\setminus\ell}
-\mathrm{Catch}_{\Pi}$ (Equation~\ref{eq:delta-catch}). Because the ablation
removes exactly one guard predicate from the conjunction of
Equation~\ref{eq:safety-invariant}, the delta isolates the marginal effect of
that predicate on the strict category-match criterion. The paper drafts report
an approximate run-to-run noise band of $\pm1$--$2$ percentage points for the
ablation configuration, which is carried into the results chapter as an
interpretation aid rather than a formal confidence statement
(Section~\ref{sec:data-analysis}).

\subsection*{Operational measures}

In addition to the accuracy metrics, the harness records operational cost for
each cache-miss request: end-to-end latency from request receipt to the returned
decision, aggregated as percentiles over the run. These figures are recorded per
backend and configuration in the evaluation manifest. Because the author has not
yet finalised the latency figures for the canonical 1500-record evaluation, no
latency percentile is quoted in this dissertation at present; the values will be
reported in Chapter~\ref{ch:results} once they are confirmed against the
manifest. Semantic-cache-hit latency is recorded separately, since a cache hit
bypasses model inference entirely.

\TODO{Latency figures for the canonical evaluation are PENDING from the author
(AUDIT checkpoint: quote them in this section and in Chapter~\ref{ch:results}
only after the author confirms the values against the evaluation manifest).}

\section{Data Analysis}\label{sec:data-analysis}

The analysis applied to the recorded measurements is descriptive and
quantitative: every claim in the results chapter is a comparison of rates
computed from the committed experiment record. In line with the design decision
of Section~\ref{sec:research-design}, no inferential statistics are applied---no
significance tests, no confidence intervals and no seed sweeps are reported.
Each experiment therefore yields point estimates, and the analysis states,
for each comparison, the rule by which a difference is read.

\subsection*{Pairwise comparisons}

Experiment~1 (raw against guarded) and Experiment~2 (Gemma against Qwen under an
identical stack) are read as side-by-side rates per backend and per metric.
Because the pipeline's stochastic component is sampled once per record and
condition, the reported differences are point comparisons; the analysis does not
attach uncertainty bounds to them. The results are presented as point deltas
(percentage points) and interpreted qualitatively where the underlying
mechanism is deterministic: for a guard whose outcome is a pure function of its
input, a difference on the affected subset is read as exact behaviour, not as a
sampling artefact.

\subsection*{Ablation deltas}

For Experiment~3, the analysis uses the per-layer delta of
Section~\ref{sec:procedure}, $\Delta\mathrm{Catch}_{\ell}$ (Equation
\ref{eq:delta-catch}), together with the accompanying movement in
$D_{\mathrm{Acc}}$ and FP. The recorded run-to-run noise band of approximately
$\pm1$--$2$ percentage points is the interpretation boundary: a delta outside the
band is read as a layer effect, while a delta inside the band is not
distinguished from noise. This is a pragmatic reading rule, not a statistical
one, and is carried into the results chapter as such. Two deltas require special
handling. First, the combined validation-stack condition disables S5--S7 as one
block, so its delta is attributed to the block as a whole, never to a single
stage. Second, a \emph{positive} delta (catch \emph{rises} when a layer is
removed) signals an interaction rather than a contribution; the analysis records
where this pattern appears---the paper runs observe it for the scope guard, the
network validator, the self-consistency gate and the semantic cache in both
backends, with FP rising alongside it---and treats those layers as
false-positive suppressors until the underlying mechanism is confirmed. The
candidate explanation (a pre-emptive generic rejection prevents the later stage
from assigning the record's true specific label) is recorded as a hypothesis,
not an established result.

\subsection*{Disaggregation}

Two disaggregations are fixed in advance because the canonical drafts report
them. Catch is analysed per flaw category against each category's own
denominator (Experiment~4), never as an aggregate number alone, because the
categories carry different denominators and the weak semantic classes would be
hidden in a single figure. The $D_{\mathrm{Acc}}$ layer ordering is analysed
\emph{per backend} rather than as one ranking: the canonical drafts find no
single layer that dominates accuracy across both backends, so a combined ranking
would mislead. Consistency claims across backends are restricted to the ordering
of catch-rate effects, and any claim about accuracy ordering is stated
backend-specifically.

\subsection*{Result-to-question mapping}

The analysis produces the evidence for each research question as follows: RQ1 is
answered by the raw-versus-guarded deltas of Experiment~1; RQ2 by the ablation
deltas, the per-class disaggregation and the recorded interactions of
Experiment~3--4; and RQ3 by comparing whether the ordering of effects carries
across the two backends (Experiment~2 and the per-backend ablation columns of
Experiment~3). The tables that instantiate this analysis are generated directly
from the committed record by the reproduction command
(Section~\ref{sec:reproducibility}); the interpretation of validity threats that
the analysis cannot address is deferred to Chapter~\ref{ch:discussion}.

\section{Ethics, Privacy, Security and Responsible AI}\label{sec:ethics}

\subsection*{Data ethics}

The study involves no human participants and processes no personal data. The
evaluation data set (Section~\ref{sec:data-materials}) is author-authored and
synthetic: every record is a constructed security request, not a sample of
operational traffic, and no record captures information about an identifiable
individual. No consent procedures, participant recruitment or data-protection
clearance for personal data were therefore required. The synthetic set is
versioned with the project source so that the same records, and only those
records, underpin every measurement reported in this dissertation
(Section~\ref{sec:reproducibility}).

\subsection*{Dual use}

The subject of the dissertation is a natural-language gateway to security
tooling, including scanning, reconnaissance and the drivers of exploitation
frameworks within an operator-declared scope. The capability is dual use by
nature. Three design and practice decisions bound this risk. First, the system
is deployment-gated: every action is subject to role-based access control and
engagement-scope enforcement, and the safety invariant
(Equation~\ref{eq:safety-invariant}) ensures that an accepted intent is
authorised and in scope before anything downstream may act on it. Second, the
evaluation is conducted entirely offline against synthetic prompts in a lab
harness; no live, third-party or production system was ever contacted, and the
design assumption that downstream components honour the rejected-intent contract
is declared in Section~\ref{sec:scope-limitations-assumptions}. Third, the
recorded failure modes---including the residual weakness of the semantic flaw
classes reported in Chapter~\ref{ch:results}---are published as part of this
dissertation rather than suppressed, in keeping with established practice in
dual-use security research.

\subsection*{Privacy}

The pipeline is deliberately offline and single-device: no utterance or result
leaves the device, nothing is retrieved from an external service at inference
time, and the operator's requests are not transmitted anywhere by the system
(Section~\ref{sec:research-design}). The only persistent assets are the model
weights, the operator-configured policies and the session context. No personal
data is collected, stored, shared or required for any experiment reported here.

\subsection*{Security in the research process}

The study is conducted under a defined threat model (Sections~\ref{sec:gap-synthesis}
and~\ref{sec:research-design}): the attacker controls the single user utterance,
not the device, prompts, policies or model weights. The evaluation harness
exercises the pipeline as a lab component, with no free-form model output ever
rendered to a downstream shell and a per-stage timeout guarding against a
stalled model. No secrets beyond the model weights are held, and the research
artefacts (code, data set, evaluation scripts and manifest) are versioned so the
experiment record can be audited independently.

\subsection*{Responsible use of artificial intelligence}

The guarded design is itself a responsible-AI measure: the model's output is
treated as an untrusted proposition until validated deterministically, which is
the stance this dissertation argues every LLM-integrated security automation
should adopt (Sections~\ref{sec:gap-synthesis} and~\ref{ch:product-design}).
Responsible reporting is observed throughout: measurements are presented as
single-run point estimates without overstated confidence
(Sections~\ref{sec:research-design} and~\ref{sec:data-analysis}), limitations
and residual failure modes are declared rather than hidden
(Section~\ref{sec:scope-limitations-assumptions}), and the safety/usability
trade-off is reported honestly, including the deliberate preference for
conservative over-blocking over risky under-blocking (Section~\ref{sec:measures}).

In preparing this dissertation, AI-based tools were used for language
correction, drafting support and the development of the artefacts described
here. The use follows the institution's policy on AI tools and is disclosed in
the applicable declaration (Appendix~\ref{app:b}).

\subsection*{Compliance}

Because the study involves no human participants, no personal data and no
intervention in live systems, it falls outside the class of work requiring
institutional research-ethics approval, and the applicable integrity declaration
is included in Appendix~\ref{app:b}.

\TODO{At the final AUDIT step, confirm with the faculty (i) that no
institutional research-ethics approval is required for a purely synthetic,
no-participant study of this kind, and (ii) the exact wording required for the
AI-tools disclosure in Appendix~\ref{app:b} under the current institutional AI
policy, and align the disclosure in this section with it.}

\section{Reproducibility and Research Artefacts}\label{sec:reproducibility}

The study is reproducible by design. Every material that produced the
measurements reported in this dissertation is versioned with the project
source, and a single command re-executes the evaluation and regenerates every
table used in the results chapter. The purpose of this section is to state
exactly what the artefacts are, what the evaluation manifest pins, and what the
boundaries of reproducibility are given the stochastic component of the system.

\subsection*{Artefact set}

The following artefacts constitute the research record:
\begin{itemize}
\item \textbf{Source code.} The Modular Python service (FastAPI) with the
separation between core stages, guard middleware and the language-model adapter
(Chapter~\ref{ch:implementation-testing}).
\item \textbf{Data set.} The $N{=}1500$ annotated evaluation data set
(Section~\ref{sec:data-materials}) together with the synthetic-generation
script used to author it.
\item \textbf{Evaluation harness.} The scripts that execute every experiment,
apply the procedure of Section~\ref{sec:procedure} and compute the metrics of
Section~\ref{sec:measures}.
\item \textbf{Evaluation manifest.} The single configuration document that pins
every experimental constant (below).
\item \textbf{Committed experiment record.} The per-record decisions written by
the harness, from which every reported table is regenerated without re-running
the model.
\item \textbf{Automated test suite.} The 482 automated tests (434 unit and 48
asynchronous integration tests, measured 2026-08-26) that exercise the guard
paths, the RBAC matrix, the cache paths and the hallucination policy, and that
gate the harness before any measurement is accepted.
\end{itemize}

\subsection*{Evaluation manifest}

The manifest is the authoritative anchor for reproducing the study. It pins the
Ollama runtime version and the per-model image digests as fetched, the hardware
platform and operating system, the data set version, the sampling temperatures
($0.4$ primary and $0.6$ for the $k{=}3$ consistency samples), the semantic
agreement threshold $\theta_s$, the adversarial-detector flag and auto-block
thresholds, the operator-configured engagement scope and RBAC policy, and the
per-stage timeout. No experimental constant is decided at run time
(Section~\ref{sec:research-design}).

\subsection*{Reproduction procedure}

A clean checkout of the versioned artefacts is sufficient: a single reproduction
command runs the harness against the data set, recomputes every metric and
regenerates every table that appears in Chapter~\ref{ch:results}. Because the
per-record decisions are committed, the reported figures can also be regenerated
directly from the experiment record without re-invoking the model, which makes
the numerical claims of this dissertation auditable independently of the
stochastic generator. The evaluation instruments used in that record are
collated in Appendix~\ref{app:a}.

\subsection*{Boundaries of reproducibility}

Two boundaries are stated plainly. First, the model component is stochastic:
an independent re-run may yield point estimates that differ slightly from the
single-run values committed here, which is precisely why the dissertation
reports point estimates without confidence claims (Sections~\ref{sec:research-design}
and~\ref{sec:data-analysis}); the committed record guarantees that the values
reported are real, measured values. The deterministic guard layers, by contrast,
are exact functions of their inputs and reproduce bit-for-bit on the same input
set. Second, absolute latency figures are bound to the evaluation hardware and
runtime; the manifest pins the exact versions so that the record can be
reproduced on comparable equipment, and cross-device comparisons require
re-measurement rather than extrapolation (Section~\ref{sec:measures}).

\subsection*{Integrity}

The research artefacts contain no personal or sensitive data, and no record in
the data set corresponds to an identifiable individual. The versioned artefact
set, the committed experiment record and the reproduction command together
constitute an auditable trail from the raw measurements to every number quoted
in this dissertation; discrepancies between the record and the text are defects
to be corrected, not adjustments to be made. With the protocol fixed in this
chapter, the following chapters develop the design of the system under test
(Chapter~\ref{ch:product-design}), its implementation (Chapter~\ref{ch:implementation-testing}),
and the measurements and their interpretation (Chapters~\ref{ch:results}
and~\ref{ch:discussion}).

\TODO{At the final AUDIT step, verify on a clean checkout that the single
reproduction command executes end-to-end, regenerates every table cited in
Chapter~\ref{ch:results}, and that the committed experiment record is complete
for every configuration reported.}